\section{IWRS / IVRS}

% SLIDE ID: sec05-goal
\BlockGoalFrame{\SlideID{sec05-goal}}{
  \item zrozumieć rolę systemów randomizacji w badaniu klinicznym,
  \item zobaczyć, jak system zarządza przypisaniem leku badanego,
  \item omówić ryzyka operacyjne i skutki błędów użytkownika.
}

% SLIDE ID: sec05-ivrs
\begin{frame}{Czym jest IVRS}
  \SlideID{sec05-ivrs}
  \begin{itemize}
    \item IVRS to telefoniczny system interaktywny używany historycznie i nadal w części badań,
    \item użytkownik wykonuje operacje przez połączenie głosowe i wybór opcji,
    \item system wspiera randomizację i logistykę leku bez konieczności pracy w interfejsie webowym.
  \end{itemize}
\end{frame}

% SLIDE ID: sec05-iwrs
\begin{frame}{Czym jest IWRS}
  \SlideID{sec05-iwrs}
  \begin{itemize}
    \item IWRS to internetowy system randomizacji i zarządzania lekiem badanym,
    \item zwykle przejmuje funkcje IVRS w nowoczesnych badaniach,
    \item łączy logikę randomizacji, zapasy leku i zapis operacji w jednym miejscu.
  \end{itemize}
\end{frame}

% SLIDE ID: sec05-history-overview
\begin{frame}{Historia rozwoju IVRS / IWRS / RTSM}
  \SlideID{sec05-history-overview}
  \begin{center}
    \Large
    Ręczne zarządzanie $\rightarrow$ telefoniczne \texttt{IVRS} $\rightarrow$ webowe \texttt{IWRS} $\rightarrow$ współczesne \texttt{RTSM / IRT}
  \end{center}

  \vspace{0.45cm}

  \begin{itemize}
    \item Historia tych systemów pokazuje, że randomizacja i logistyka leku nie mogły pozostać procesami ręcznymi przy rosnącej złożoności badań.
    \item Z czasem ciężar systemu przesunął się z samego przypisania leczenia na szersze zarządzanie zapasem, resupply i kontrolę nad całym przepływem leku badanego.
    \item Dlatego współczesne \texttt{IWRS} coraz częściej są opisywane szerzej jako \texttt{RTSM} albo \texttt{IRT}.
  \end{itemize}
\end{frame}

% SLIDE ID: sec05-history-pre-digital
\begin{frame}{1. Era przedcyfrowa: przed latami 90.}
  \SlideID{sec05-history-pre-digital}
  \begin{itemize}
    \item Przed automatyzacją personel ośrodka i zespoły operacyjne zarządzały randomizacją oraz lekiem badanym ręcznie.
    \item Zestawy leku były wysyłane do ośrodków z góry, często w pełnych blokach odpowiadających planowanemu leczeniu pacjentów.
    \item Jeśli ośrodek rekrutował słabo albo pacjenci odpadali z badania, rosło marnotrawstwo zapasu i ryzyko nieefektywnej dystrybucji.
    \item Ten model był mało elastyczny, słabo skalowalny i utrudniał utrzymanie bieżącej kontroli nad dostępnością leku.
  \end{itemize}
\end{frame}

% SLIDE ID: sec05-history-ivrs
\begin{frame}{2. Wzrost znaczenia IVRS: lata 90.}
  \SlideID{sec05-history-ivrs}
  \begin{itemize}
    \item W latach 90. pojawiły się pierwsze systemy \texttt{IVRS}, oparte na połączeniach telefonicznych i wyborach z klawiatury numerycznej.
    \item Użytkownik dzwonił na numer systemu, wprowadzał dane pacjenta i otrzymywał przydział leczenia lub numer zestawu w czasie rzeczywistym.
    \item Najważniejszą korzyścią była automatyzacja randomizacji bez ujawniania przydziału leczenia osobom, które nie powinny go znać.
    \item Z czasem \texttt{IVRS} zaczęło wspierać nie tylko randomizację, ale także dispensing i podstawowy resupply na poziomie ośrodka.
  \end{itemize}
\end{frame}

% SLIDE ID: sec05-history-iwrs
\begin{frame}{3. Rewolucja webowa: od IWRS do RTSM}
  \SlideID{sec05-history-iwrs}
  \begin{itemize}
    \item Wraz z upowszechnieniem internetu pojawiły się webowe interfejsy \texttt{IWRS}, które początkowo bywały dodatkiem lub backupem dla rozwiązań telefonicznych.
    \item Z czasem to właśnie interfejs webowy stał się dominującym standardem pracy operacyjnej w badaniach klinicznych.
    \item Kluczowa zmiana dotyczyła logistyki: system przestał tylko losować pacjenta i zaczął aktywnie sterować zapasem, resupply i ruchem leku między depot a ośrodkiem.
    \item Dlatego coraz częściej mówi się o \texttt{RTSM} (\texttt{Randomization and Trial Supply Management}) albo szerzej o \texttt{IRT} jako parasolowym terminie dla tych technologii.
  \end{itemize}
\end{frame}

% SLIDE ID: sec05-history-modern
\begin{frame}{4. Współczesne systemy IWRS / RTSM / IRT}
  \SlideID{sec05-history-modern}
  \begin{itemize}
    \item Dzisiejsze systemy wspierają nie tylko randomizację i przydział leku, ale również forecasting, zarządzanie zapasem, expiry management i ścieżki wyjątków.
    \item Integrują się z \texttt{EDC}, depotami, systemami dostaw i innymi elementami ekosystemu \texttt{eClinical}.
    \item W bardziej złożonych badaniach mogą obsługiwać adaptacyjne schematy randomizacji, kontrolę temperatury, wymiany zestawów i złożoną logistykę biologiczną.
    \item W praktyce są dziś technologicznymi hubami łączącymi decyzję randomizacyjną z realnym przepływem leku badanego.
  \end{itemize}
\end{frame}

% SLIDE ID: sec05-history-what-changed
\begin{frame}{Dlaczego ta ewolucja ma znaczenie}
  \SlideID{sec05-history-what-changed}
  \begin{itemize}
    \item Zmieniła precyzję i bezstronność randomizacji, bo logika systemu zastąpiła proces ręczny.
    \item Zmieniła ekonomię badania, bo lek nie musi być wysyłany „na zapas” w takim stopniu jak dawniej.
    \item Zmieniła bezpieczeństwo operacyjne, bo łatwiej śledzić wydanie, zwrot, zniszczenie i odchylenia od planu.
    \item Zmieniła też zakres odpowiedzialności: nowoczesny \texttt{IWRS} to nie tylko narzędzie do przydziału ramienia, ale centrum sterowania supply chain badania.
  \end{itemize}
\end{frame}

% SLIDE ID: sec05-randomization
\begin{frame}{Randomizacja, stratyfikacja, zaślepienie}
  \SlideID{sec05-randomization}
  \begin{itemize}
    \item randomizacja przypisuje pacjenta do ramienia badania według ustalonych reguł,
    \item stratyfikacja pomaga utrzymać równowagę między grupami,
    \item zaślepienie chroni badanie przed ujawnieniem przydziału leczenia.
  \end{itemize}
\end{frame}

% SLIDE ID: sec05-drug-management
\begin{frame}{Zarządzanie lekiem badanym}
  \SlideID{sec05-drug-management}
  \begin{itemize}
    \item system przypisuje numer zestawu lub numer opakowania,
    \item kontroluje dostępność zapasu w ośrodku,
    \item pozwala śledzić wydanie, zwrot i status leku badanego.
  \end{itemize}
\end{frame}

% SLIDE ID: sec05-unblinding
\begin{frame}{Emergency unblinding}
  \SlideID{sec05-unblinding}
  \begin{itemize}
    \item to kontrolowane ujawnienie przydziału leczenia w sytuacji medycznie uzasadnionej,
    \item nie jest zwykłą operacją administracyjną,
    \item każda taka akcja musi być uzasadniona, zapisana i możliwa do audytu.
  \end{itemize}
\end{frame}

% SLIDE ID: sec05-demo-flow
\begin{frame}{Demo flow: od kwalifikacji do przypisania leku}
  \SlideID{sec05-demo-flow}
  \begin{columns}[T,onlytextwidth]
    \column{0.18\textwidth}
    \begin{beamercolorbox}[wd=\linewidth,sep=0.6em,rounded=true]{block body}
      \centering
      \textbf{Kryteria}
    \end{beamercolorbox}

    \column{0.18\textwidth}
    \begin{beamercolorbox}[wd=\linewidth,sep=0.6em,rounded=true]{block body}
      \centering
      \textbf{Randomizacja}
    \end{beamercolorbox}

    \column{0.18\textwidth}
    \begin{beamercolorbox}[wd=\linewidth,sep=0.6em,rounded=true]{block body}
      \centering
      \textbf{Ramię}
    \end{beamercolorbox}

    \column{0.18\textwidth}
    \begin{beamercolorbox}[wd=\linewidth,sep=0.6em,rounded=true]{block body}
      \centering
      \textbf{Lek}
    \end{beamercolorbox}

    \column{0.18\textwidth}
    \begin{beamercolorbox}[wd=\linewidth,sep=0.6em,rounded=true]{block body}
      \centering
      \textbf{Zapis}
    \end{beamercolorbox}
  \end{columns}

  \vspace{0.45cm}

  \begin{itemize}
    \item pacjent spełnia kryteria,
    \item system uruchamia randomizację,
    \item przypisywane jest ramię badania i numer zestawu,
    \item cała operacja zostaje zapisana w systemie.
  \end{itemize}
\end{frame}

% SLIDE ID: sec05-operational-risks
\begin{frame}{Ryzyka operacyjne}
  \SlideID{sec05-operational-risks}
  \begin{itemize}
    \item randomizacja niezgodna z protokołem,
    \item przypisanie niewłaściwego zestawu leku,
    \item błędny zwrot lub odpisanie zapasu,
    \item nieuprawnione ujawnienie przydziału leczenia.
  \end{itemize}
\end{frame}

% SLIDE ID: sec05-consequences
\begin{frame}{Konsekwencje błędów}
  \SlideID{sec05-consequences}
  \begin{itemize}
    \item zagrożenie dla bezpieczeństwa pacjenta,
    \item naruszenie zaślepienia,
    \item problemy z rozliczeniem leku badanego,
    \item ryzyko odchyleń protokołowych i wpływu na wiarygodność badania.
  \end{itemize}
\end{frame}

% SLIDE ID: sec05-demo
\DemoFrame{Randomizacja i przypisanie leku}{
  kwalifikacja pacjenta, uruchomienie randomizacji, przypisanie ramienia badania, nadanie numeru zestawu oraz zapis całej operacji w systemie.
}

% SLIDE ID: sec05-screens
\begin{frame}{Jakie widoki IWRS / RTSM warto pokazać}
  \SlideID{sec05-screens}
  \begin{itemize}
    \item ekran kwalifikacji i uruchomienia randomizacji,
    \item ekran przypisania ramienia badania i numeru zestawu,
    \item widok zapasu leku w ośrodku i statusu resupply,
    \item raport wydania, zwrotu i reconciliation leku,
    \item ekran \texttt{emergency unblinding} lub jego kontrolowanej ścieżki,
    \item dashboard pokazujący połączenie randomizacji z logistyką dostaw.
  \end{itemize}
\end{frame}

% SLIDE ID: sec05-case
\CaseStudyFrame{Co się stanie, jeśli użytkownik omyłkowo odeśle lek albo ujawni randomizację?}{
  Użytkownik wykonuje niewłaściwą operację logistyczną albo poznaje przypisanie pacjenta do ramienia badania poza kontrolowaną procedurą.
}{
  Może dojść do zaburzenia zapasu, błędnego wydania leku, naruszenia zaślepienia i konieczności formalnego wyjaśnienia incydentu w dokumentacji badania.
}

% SLIDE ID: sec05-sources
\begin{frame}{Źródła do tej sekcji}
  \SlideID{sec05-sources}
  \footnotesize
  \begin{itemize}
    \item 4G Clinical, \textit{IRT. IVRS. IXRS. IWRS. RTSM. What? Why do we do this to ourselves?}, 5 września 2017.
    \item Clinical Leader / Synteract, \textit{IVRS/IWRS Services}.
    \item Clinical Leader / Aptuit, \textit{Fully Web-Integrated Interactive Voice Response System (IVRS/IWRS)}, 2011.
    \item Clinical Leader / Cenduit, \textit{Drug Supply Management and Allocation}.
  \end{itemize}
\end{frame}
